\documentclass[11pt]{article}
\usepackage[margin=1in]{geometry}
\usepackage{amsmath,amssymb,amsthm,booktabs}
\usepackage[hidelinks]{hyperref}
\newtheorem{proposition}{Proposition}
\newtheorem{corollary}{Corollary}
\title{A six-cycle obstruction for 2-factor Hamiltonicity in cubic bipartite braces}
\author{Claude Make No Mistakes}
\date{25 November 2026}
\begin{document}
\maketitle

\begin{abstract}
We characterize when a six-cycle in a cubic bipartite brace of girth
at least six cannot be a component of a 2-factor. Its exterior must
split into two connected shores with prescribed six-edge cuts.
We use this decomposition to search at order 46 and certify a
disconnected 2-factor for each of 26,010 sampled labeled graphs.
Among these graphs are 249 certified braces with no six-cycle in
any 2-factor.
\end{abstract}

\section{Setting}
A brace here is a connected bipartite graph in which every matching
of at most two edges extends to a perfect matching. The
Funk--Jackson--Labbate--Sheehan conjecture \cite{fjls} is equivalent to
asserting that \(K_{3,3}\) and the Heawood graph are the only cubic
bipartite braces in which all 2-factors are Hamiltonian \cite{gjw}.
Abreu et al.\ verified the conjecture through order 42, and through
order 52 for girth at least eight \cite{abreu}.
The order restriction \(2\bmod 4\) for a 2-factor Hamiltonian cubic
bipartite graph \cite{fjls,gjw} singles out order 46 at girth six as
the first case outside those ranges.

\begin{proposition}
Let \(G\) be a simple cubic bipartite brace of girth at least six, and
let \(C\) be a six-cycle. The cycle \(C\) fails to occur as a component
of a 2-factor if and only if \(V(G)\setminus V(C)\) can be partitioned
into nonempty connected shores \(X,Y\) with the following properties:
\begin{enumerate}
\item Each shore has three edges to \(C\) and three to the other shore.
\item All boundary endpoints within \(X\) have one bipartition color,
and all boundary endpoints within \(Y\) have the opposite color.
\end{enumerate}
\end{proposition}

\begin{proof}
Write \(A,B\) for the bipartition, with \(\lvert A\rvert=\lvert B\rvert=m\).
We first recall the brace expansion
\(\lvert N(Q)\rvert\geq \lvert Q\rvert+2\) for any nonempty set \(Q\)
in one color class with \(\lvert Q\rvert\leq m-2\).
Cubic degree counting gives \(\lvert N(Q)\rvert\geq\lvert Q\rvert\).
Equality would disconnect \(G\). If \(\lvert N(Q)\rvert=\lvert Q\rvert+1\),
there are three edges from \(N(Q)\) to the exterior of \(Q\).
They include two disjoint edges. Otherwise the three share an
endpoint. A common endpoint in \(N(Q)\) would have no neighbor in \(Q\).
A common endpoint outside \(Q\) would make \(Q\cup N(Q)\) together with
that endpoint a proper component. Extending two disjoint crossing edges
to a perfect matching leaves too few neighbors to cover \(Q\).

By girth, \(C\) is chordless and its six external neighbors are
distinct. Write \(D_A,D_B\) for the external neighbors by color.
A perfect matching complementary to a 2-factor containing \(C\)
must contain the six edges leaving \(C\). These edges extend
if and only if \(H=G-(V(C)\cup D_A\cup D_B)\) has a perfect matching.
The remaining color classes \(A_H,B_H\) have size \(m-6\).

\newpage
If the matching fails, Hall's theorem gives a nonempty
\(S\subseteq A_H\) with \(\lvert N_H(S)\rvert<\lvert S\rvert\).
The only neighbors \(S\) can lose are the three in \(D_B\).
Brace expansion forces \(\lvert N_H(S)\rvert=\lvert S\rvert-1\)
and \(D_B\subseteq N_G(S)\).
Put \(T=B_H\setminus N_H(S)\). Then
\(\lvert T\rvert=m-6-\lvert S\rvert+1\) and
\(N_H(T)\subseteq A_H\setminus S\). Applying the same argument
to \(T\) gives \(N_H(T)=A_H\setminus S\) and \(D_A\subseteq N_G(T)\).
Thus \(X=S\cup N_G(S)\) and \(Y=T\cup N_G(T)\) partition the exterior.
Only vertices of \(N_G(S)\) can meet the cut of \(X\).
The color imbalance \(\lvert N_G(S)\rvert-\lvert S\rvert=2\)
and cubic degrees give exactly six cut edges. Three go to \(C\)
through \(D_B\). Symmetrically \(Y\) has six cut edges, three
through \(D_A\).

Every component of either shore has a positive cut size divisible
by three, by the same color and degree count. If a shore is
disconnected it has two components, each with three cut edges.
A nonsingleton component would have color imbalance one
and would contradict brace expansion on its smaller color class.
Two singleton components are impossible since each shore
has at least four vertices.

Conversely, the given boundary conditions imply color imbalance two
on each shore. A perfect matching has to use exactly two edges of
each of these cuts. The three edges from \(C\) entering either shore
already exceed that number. Therefore the six edges leaving \(C\)
cannot extend.
\end{proof}

One can test \(C\) using either alternating triple of its six
external edges. If a perfect matching contains the triple at the
three cycle vertices of one color, its complement contains both
cycle edges at each of those vertices and hence all of \(C\).

\begin{corollary}
For a blocked six-cycle, let the shore orders be \(2s+2,2t+2\) with
\(s,t\) the smaller color-class sizes. Each of \(s,t\) is either
1 or at least 4. A shore with \(s=1\) is a claw. A shore with
\(s=4\) is the graph obtained by subdividing every edge of \(K_4\) once.
If \(\lvert V(G)\rvert>14\), there is a cyclic six-edge cut with
both sides of order at least 10.
\end{corollary}

\begin{proof}
The smaller class in a shore has every neighbor inside that shore.
Its vertices have neighborhoods of size three, with pairwise
intersections at most one since \(G\) has girth at least six.
Thus \(s+2\geq 3s-\binom{s}{2}\), which excludes \(s=2,3\).
For \(s=4\), equality forces each of the six pairs to share
a neighbor and every such neighbor to belong to exactly two of
the neighborhoods. This gives the subdivision claim. If the
order is greater than 14, some shore has at least ten vertices.
Its six-edge cut leaves at least ten vertices on the other side,
which contains \(C\). A shore of \(k\geq10\) vertices in a cubic
graph with six cut edges has \((3k-6)/2\) internal edges and
contains a cycle.
\end{proof}

At order 46, \(s+t=18\). Up to interchanging the shores, the smaller
parameter belongs to \(\{1,4,5,6,7,8,9\}\). The larger shore has at
least 20 vertices.

\newpage
\section{Search and exact checks}
We sample bipartite adjacency lists for the two shores and join them
using the six-cycle and the prescribed edges. We discard invalid,
disconnected, or girth-four graphs. From low-score examples we
apply bipartition-preserving two-edge switches that keep the
defining six-cycle and its shore decomposition. The score counts
six-cycles that occur as a component of a 2-factor. In the final
run we visit switch neighborhoods of graphs with score at most one,
subject to a cap on the number of new graphs.

\begin{center}
\begin{tabular}{lrrr}
\toprule
Stage & Graphs & Six-cycle witness & Longer-cycle witness \\
\midrule
Direct sampling & 763 & 763 & 0 \\
Switch walks & 7,247 & 7,247 & 0 \\
Switch neighborhoods & 18,000 & 17,751 & 249 \\
\bottomrule
\end{tabular}
\end{center}

Each rejection includes a perfect matching whose complement is
disconnected. A separate verifier checks every witness. For
the 249 zero-score graphs, it independently enumerates
six-cycles and tests the matching extension of an alternating
triple by breadth-first augmentation. It tests 2,208 disjoint edge
pairs in each zero-score graph to verify brace status. Each of
these graphs has 12 six-cycles. The separately archived example
has a 2-factor with cycle orders 22 and 24. The sampled graphs are
distinguished by their labeled adjacency lists. The search does
not exhaust order 46.

The proof uses the complementary Hall-set and degree-count
mechanism also seen in generalized matching-extension results
\cite{fks}. The criterion here describes the exact obstruction
for a fixed six-cycle in a cubic bipartite brace.

\begin{thebibliography}{9}
\bibitem{fjls} M. Funk, B. Jackson, D. Labbate, and J. Sheehan.
2-factor Hamiltonian graphs. \emph{J. Combin. Theory Ser. B}
87 (2003), 138--144.
\href{https://doi.org/10.1016/S0095-8956(02)00031-X}{DOI link}.
\bibitem{gjw} M. Gorsky, T. Johanni, and S. Wiederrecht.
A note on the 2-factor Hamiltonicity Conjecture.
\emph{Discrete Math.} 348 (2025), 114442.
\href{https://arxiv.org/html/2408.08128v2}{Preprint}.
\bibitem{abreu} M. Abreu, J. Goedgebeur, J. Jooken, F. Romaniello,
and T. Van den Eede. The Gray graph is pseudo 2-factor isomorphic.
\href{https://arxiv.org/html/2504.12095v2}{arXiv:2504.12095v2}, 2026.
\bibitem{fks} J. Fujisawa, S. Koizumi, and Y. Suzuki.
Generalized matching extension problem in regular bipartite graphs.
\emph{Discrete Math.} 349 (2026), 115182.
\href{https://doi.org/10.1016/j.disc.2026.115182}{DOI link}.
\end{thebibliography}
\end{document}
